\def\topic{}
\ifdefined\MyHeaderIncluded
\else
\def\qwertyflag{true}

\ifdefined\MyHeaderIncluded

\else
  \newcommand{\foobar}{FooBar}
  \def\MyHeaderIncluded{true}
% if is till the end of the file

\documentclass[10pt,a4paper]{article}
\usepackage[utf8]{inputenc}
%\usepackage[german]{babel}
%\usepackage[cyrillic]{babel}
\usepackage[T1]{fontenc}
%\usepackage[T2A]{fontenc}
%\usepackage[english, russian]{babel}
\usepackage{amsmath}
\usepackage{amsfonts}
\usepackage{amssymb}
\usepackage{amsthm}
\usepackage{mathtools} %coloneqq

\usepackage{bussproofs} %

%\renewcommand{\thesubsection}{\arabic{subsection}}
%\newtheorem{thm}{Theorem}[section]
%\renewcommand{\thethm}{\arabic{subsection}.\arabic{thm}}

\usepackage{thmtools}
\usepackage{hyperref}
\hypersetup{
    colorlinks=true,
    linkcolor=blue,
    filecolor=magenta,      
    urlcolor=cyan,
    pdftitle={Overleaf Example},
    %pdfpagemode=FullScreen,
    }
%\usepackage{stix-type1}%\turnediota
\theoremstyle{definition}

\newtheorem{constructor}{Constructor}[section]
\newtheorem{iaxiom}{Axiom}
\newtheorem{axiom}{Axiom}
\newtheorem*{axiom*}{Axiom}
\newtheorem{definition}{Definition}[section]
\newtheorem*{definition*}{Definition}
\newtheorem{remark}{Remark}[section]
\newtheorem*{remark*}{Remark}
\newtheorem{recall}{Recall}[section]
\newtheorem*{recall*}{Recall}

\newtheorem{lemma}{Lemma}[section]
\newtheorem{theorem}{Theorem}[section]
\newtheorem*{theorem*}{Theorem}
\newtheorem{note}{Note}[section]
\usepackage{etoolbox}

%\newcommand*{\@Dashv}[2]{%
%  \reflectbox{$\m@th#1#2$}%
%}
\newcommand*{\separation}[3]{% 1)binder, 2)set, 3)condition
  \{#1\in #2: #3\}
}

% https://www.freecodecamp.org/news/simple-regex-tricks-for-beginners-3acb3fa257cb/
%  \\mathcal\{T\} |---> \topology
%  \\mathcal\{([^\}]*)\}   |--->  $1
\newcommand{\AND}{\ \ \&\ \ }
\newcommand{\any}{\mathrm{\ any}}
\newcommand{\borel}{\mathcal{B}}
\newcommand{\can}{\mathrm{can}}
\newcommand{\canonize}{\mathrm{canonize}}
\newcommand{\closure}{\mbox{cl}}
\newcommand{\complex}{\mathbb{C}}
\newcommand{\converges}{\longrightarrow}
\newcommand{\countable}{\mbox{\ countable}}
\newcommand{\ldefined}{\mbox{\ defined}}
\newcommand{\dom}{\mbox{dom}}
\newcommand{\Dom}{\mbox{Dom}}
\newcommand{\Done}{\noindent \textbf{Done.}\\}
\newcommand{\lexists}{\mbox{\ exists}}
\newcommand{\lforall}{\mbox{\ forall}}
\newcommand{\lunique}{\mbox{\ unique}}
\newcommand{\family}{\mathrm{\ family}}
\newcommand{\Fm}{\mathrm{Fm}}
\newcommand{\FuncSymb}{\mathrm{FuncSymb}}
\newcommand{\function}{\mathrm{\ function}}
\newcommand{\functional}{\mbox{\ functional}}
\newcommand{\IFF}{\ \Longleftrightarrow\ }
\renewcommand{\iff}{\ \longleftrightarrow\ }
\newcommand{\IMPLIES}{\ \Longrightarrow\ }
\renewcommand{\implies}{\ \longrightarrow\ }
\newcommand{\Inf}{\mathrm{Inf}}
\newcommand{\injective}{\mbox{\ injective}}
\newcommand{\injection}{\mbox{\ injection}}
\newcommand{\interior}{\mbox{int}}
\newcommand{\intersection}{\bigcap}
\newcommand{\isempty}{\mathrm{\ empty}}
\newcommand{\LB}{\mathrm{LB}}
\newcommand{\nghbset}{\mathcal{N}}
\newcommand{\nonempty}{\mathrm{\ nonempty}}
\newcommand{\Ord}{\mathrm{Ord}}
\newcommand{\power}{\mathcal{P}}
\newcommand{\pr}{\mathrm{pr}}
\newcommand{\fst}{\mathrm{fst}}
\newcommand{\snd}{\mathrm{snd}}
\newcommand{\PredSymb}{\mathrm{PredSymb}}
\newcommand{\Proof}{\noindent \textbf{Proof.}\\}
\newcommand{\ran}{\mathrm{ran}}
\newcommand{\Ran}{\mathrm{Ran}}
\newcommand{\real}{\mathbb{R}}
\newcommand{\relation}{\mathrm{\ relation}}
\newcommand{\set}{\mathrm{\ set}}
\newcommand{\singleton}{\mathrm{\ singleton}}
\newcommand{\semantics}{\mathrm{semantics}}
\newcommand{\semialgebra}{\mathcal{S}}
\newcommand{\Sup}{\mathrm{Sup}}
\newcommand{\surjective}{\mbox{\ surjective\ on\ }}
\newcommand{\Term}{\mathrm{Term}}
\newcommand{\Type}{\mathrm{Type}}
\newcommand{\topology}{\mathcal{T}}
\newcommand{\total}{\mbox{\ total\ on\ }}
\newcommand{\transitive}{\mathrm{\ transitive}}
\newcommand{\union}{\bigcup}
\newcommand{\lTrue}{\mathrm{True}}
\newcommand{\lFalse}{\mathrm{False}}
\newcommand{\lUncertain}{\mathrm{Uncertain}}
\newcommand{\lif}{\mathrm{if}}
\newcommand{\lthen}{\mathrm{then}}
\newcommand{\lelse}{\mathrm{else}}
\newcommand{\lexception}{\mathrm{exception}}
\newcommand{\lundef}{\mathrm{undef}}
\newcommand{\univocal}{\mathrm{\ univocal}}
\newcommand{\univocalat}{\mathrm{\ univocal\ at\ }}
\newcommand{\UB}{\mathrm{UB}}
\newcommand{\Un}{\mathcal{U}}
\newcommand{\Var}{\mathrm{Var}}

\newcommand{\notvarphi}{\phi}
\renewcommand{\phi}{\varphi}

\newcommand{\comment}[1]{\quad\mbox{(#1)}}
% local commands
%\preto\section{\undef\blipo} %etoolbox
\preto\subsection{\undef\blipo} %etoolbox
% https://tex.stackexchange.com/questions/105684/limit-macro-newcommand-scope-to-section
\usepackage{graphicx}
\usepackage{authblk}

\setlength\parindent{0pt} % whole doc of \noindent
\usepackage{soul}
%\usepackage{mathgl}

\fi % guard
\ifdefined\topic
\else
\def\topic{No Titel}
\fi
\title{\topic}
\author[1]{Georgy Dunaev}
\begin{document}
\maketitle
\fi
\section{Ordered pair}
the following definition is too weak:
\begin{definition}
\label{def:ordered pair} unsatisfactory
Ordered pair. %maybe V somehow?
\[
\langle A,B \rangle=\{x|(A\set)\land(B\set)\land (x\in\{\{A\},\{A,B\}\})\}
\]
\end{definition}

\subsection{Definitions of the ordered pair for sets and classes}
Separate definitions:
\[
(a,b) := \{(0,a_1)|a_1\in a\}\cup \{(1,b_1)|b_1\in b\}
\]
\[
X\times Y := \{(x,y)|x\in X \land y\in Y\}
\]

\noindent Simultaneous defintion:
\[
(a,b) := ({0}\times a)\cup ({1}\times b)
\]
\[
X\times Y := \{(x,y)|x\in X \land y\in Y\}
\]

\begin{remark}\ \\
Do not confuse any of previous definitions with a family $F:2\to\mathcal{U}$.
\[
\tilde{K}:= \{\{k\}| k\in K\}
\]
\[
[a,b] := ({0}\times \tilde{a})\cup ({1}\times \tilde{b})
\]
Generally speaking,
\[
(a,b) :\neq [a,b]
\]
because then there is gurantee that it is well-founded.(todo: check it)\\
\end{remark}


%Remark 1:
%\[
%(a,b) := \{(0,a_1)|a_1\in a\}\cup \{(1,b_1)|b_1\in b\}
%\]
%\[
%   = \{\{(1,a_2)|a_2\in a_1\}|a_1\in a\}\cup \{(1,b_1)|b_1\in b\}
%(0,b) = \{(0,a_1')|a_1'\in \emptyset\}\cup \{(1,b_1)|b_1\in b\} = \{(1,b_1)|%b_1\in b\}

%(1,b_1)= \{(0,a_1')|a_1'\in \{\emptyset\}\}\cup \{(1,b_2)|b_2\in b_1\}
%       = \{(0,\emptyset)\}\cup \{(1,b_2)|b_2\in b_1\}
%       = \{(0,\emptyset)\}\cup \{\{(0,\emptyset)\}\cup \{(1,b_2)|b_2\in b_1\}|%b_2\in b_1\}
%\]
\[
(a,0) = \{(0,a_1)|a_1\in a\}
\]
\[
(0,b) = \{(1,b_1)|b_1\in b\}
\]

%\Proof

Well-foundedness by epsilon-relation:\\
%Checking that $(a,b)$ is well-found.
\begin{theorem}
If $a$ and $b$ are well-founded, then $(a,b)$ is well-founded.
\end{theorem}
\Proof
We apply the epsilon-induction for $a$.\\ % maybe even  complete? nah, it's enough
Assume there is no infinite $\ni$-chain in $(a_1,b)$ for any element $a_1$ of $a$.\\
We apply the epsilon-induction for $b$.\\
Assume there is no infinite $\ni$-chain in $(a_1,b_1)$ for any element $b_1$ of $b$.\\
\[
w\in (a,b)
\]
WLOG("Case I"), 
\[
w\in \{(0,a_1)|a_1\in a\}
\]
\[
\exists a_1\in a: w\in (0,a_1)
\]

%Case 1: a is empty.\\
%Then it contradicts with $w\in \{(0,a_1)|a_1\in a\}$, so statement is exfalso proven.\\
%Case 2: a is non-empty.\\

We know that
\[
(a,\emptyset) = \{(0,a_1)|a_1\in a\}\cup \{(1,b_1)|b_1\in \emptyset\} = \{(0,a_1)|a_1\in a\}
\]
and
\[
(\emptyset,b) = \{(0,a_1)|a_1\in \emptyset\}\cup \{(1,b_1)|b_1\in b\} = \{(1,b_1)|b_1\in b\}
\]
.

Small remark: we can compute, just for the sanity check:
\[
(0,0) = (\emptyset,\emptyset) = \{(1,b_1)|b_1\in \emptyset\} = \emptyset
\]
\[
(1,0) = (\{\emptyset\},\emptyset) = \{(0,\emptyset)\} = \{\emptyset\}
\]
\[
(0,1) = (\emptyset,\{\emptyset\}) =/\mbox{by the 2nd eq}/= \{(1,\emptyset)\} = \{\{\emptyset\}\},
\]
indeed they are all different when arguments are different, as intended. In the later theorem it will be proven for arbitrary pairs.

Every chain ends with the empty set(or with an urelement in some non-extensional theories).\\

If(WA) $(0,a_1)$ has an infinite chain, then (using the second equality) 
$\{(1,e)|e\in a_1\}$ has it too. 

So there exists $e_0\in a_1$, such that $(1,e_0)$ has an infinite chain.
\[
(1,e_0) = \{(0,0)\}\cup \{(1,e_1)|e_1\in e_0\}
\]
.
Obviously, $\{(0,0)\}$ has no infinite chain. So there exists $e_1\in e_0$, such that $(1,e_1)$ has an infinite chain.
\[
(1,e_1) = \{(0,0)\}\cup \{(1,e_2)|e_2\in e_1\}
\]
and so on: (using here a natural number recursion/induction)
Obviously, $\{(0,0)\}$ has no infinite chain. So there exists $e_{n+1}\in e_{n}$, such that $(1,e_{n+1})$ has an infinite chain.
\[
(1,e_{n+1}) = \{(0,0)\}\cup \{(1,e_{n+1})|e_{n+1}\in e_n\}
\]

Therefore $a$ has an infinite chain starting with $e$, contradiction.

The proof is similar, for the "Case II", where
\[
w\in \{(1,b_1)|b_1\in b\}
\]
\[
\exists b_1\in b: w\in (1,b_1),
\]
with obtaining a contradiction with $b$ being well-founded.\\
\Done

\begin{definition}(we define a functional symbols which take the first and the second element of the pair) [we can add, "if it's an ordered pair"]
\[
\pr_1(P) := \fst(P) := \{b|(0,b)\in P\}
\]
\[
\pr_2(P) := \snd(P) := \{b|(1,b)\in P\}
\]
\end{definition}

\begin{theorem}
\[
((A_1,B_1) = (A_2,B_2)) \iff ((A_1=A_2)\land(B_1=B_2)) 
\]
\end{theorem}
\Proof
...\\
\Done

\begin{theorem}(OLD)
If $a$ and $b$ are well-founded, then $(a,b)$ is well-founded.
\end{theorem}
\Proof
\[
h((a,b)) = h(({0}\times a)\cup ({1}\times b)) = 
\]
\[
=\max( \sup\{h(y)|y\in ({0}\times a)\}, \sup\{h(y)|y\in ({0}\times a)\} ) = /\mbox{WLOG}/
\]
\[
= \sup\{h(y)|y\in ({0}\times a)\} = /\mbox{by def}/
\]
\[
= \sup\{h(y)|y\in \{(0,a_1)|a_1\in a_0\}\} =
\]
\[
= \sup\{h((0,a_1))|a_1\in a_0)\} =
\]

\Done
\ifdefined\qwertyflag
\undef\qwertyflag
\end{document}
\fi